% \iffalse meta-comment
%
% File: LuaCJK.dtx
% -----------------------------------------------------------------------
%   Copyright (C) 2026 by Mingyu Xia <myhsia@outlook.com>               *
% -----------------------------------------------------------------------
%   This work may be distributed and/or modified under the conditions   *
%   of the LaTeX Project Public License (LPPL), either version 1.3c of  *
%   this license or (at your option) any later version.                 *
%   The latest version of this license is in                            *
%                                                                       *
%       http://www.latex-project.org/lppl.txt                           *
%                                                                       *
%   and version 1.3c or later is part of all distributions of LaTeX     *
%   version 2008 or later.                                              *
%                                                                       *
%   This work has the LPPL maintenance status `maintained'.             *
%                                                                       *
%   The Current Maintainer of this work is Mingyu Xia.                  *
% -----------------------------------------------------------------------
%   This work consists of the files LuaCJK.dtx,                         *
%                                   LuaCJK.ins,                         *
%                 the derived files LuaCJK.sty,                         *
%                                   LuaCJKfntef.sty,                    *
%                                   LuaCJK.lua,                         *
%                                   LuaCJK-core.lua,                    *
%           the documentation files LuaCJK.pdf,                         *
%                               and README.md,                          *
%                               and LICENSE.                            *
% -----------------------------------------------------------------------
%
%   Any modification of this file should ensure that the copyright and  *
%   license information is placed in the derived files.                 *
%                                                                       *
% -----------------------------------------------------------------------
%
%<*internal>
\iffalse
%</internal>
%
%<*readme>
[![CTAN Version](https://img.shields.io/ctan/v/luacjk)](https://ctan.org/pkg/luacjk)
[![GitHub Release](https://img.shields.io/github/v/release/myhsia/LuaCJK)](https://github.com/myhsia/LuaCJK/releases/latest)
[![GitHub Last Commit](https://img.shields.io/github/last-commit/myhsia/LuaCJK)](https://github.com/myhsia/LuaCJK/commits)
[![Actions Status](https://github.com/myhsia/LuaCJK/workflows/Automated%20testing/badge.svg)](https://github.com/myhsia/LuaCJK/actions)
[![GitHub Repo stars](https://img.shields.io/github/stars/myhsia/LuaCJK)](https://github.com/myhsia/LuaCJK)

# LuaCJK

LuaCJK is a LuaTeX package for Chinese typesetting, inspired by `xeCJK`.

The project is currently at an early development stage. The first public
release provides a CJK font family layer, basic CJK/CJK and CJK/Latin
spacing, the `quanjiao`, `plain`, `banjiao`, `kaiming`, `CCT`, and
`hangmobanjiao` punctuation styles, and basic punctuation line-breaking
prohibitions on top of LuaTeX node callbacks.  The companion package
`LuaCJKfntef` adds underlining, emphasis dots, and strike-through through
`lua-ul`.  Fallback fonts and the remaining decoration types are planned for
later releases.

## Loading

```tex
\usepackage{LuaCJK}
\setCJKmainfont{FandolSong-Regular.otf}
```

LuaCJK requires LuaLaTeX. It must not be loaded together with `luatexja`.

## License

This work may be distributed and/or modified under the conditions of the
LaTeX Project Public License, version 1.3c or later.
%</readme>
%
%<*package|fntef>
\NeedsTeXFormat{LaTeX2e}[2026/06/01]
\RequirePackage{expl3}
\GetIdInfo $Id: LuaCJK.dtx 0.1.0 2026-11-02 myhsia<myhsia@outlook.com>$
%<package>  {A LuaTeX CJK typesetting package}
%<package>\ProvidesExplPackage {\ExplFileName}
%<fntef>  {Font effects for LuaCJK (CTEX)}
%<fntef>\ProvidesExplPackage {LuaCJKfntef}
  {\ExplFileDate} {\ExplFileVersion} {\ExplFileDescription}
%<package>\edef \DocFileDate        {\ExplFileDate}
%<package>\edef \DocFileVersion     {\ExplFileVersion}
%<package>\edef \DocFileDescription {\ExplFileDescription}
%<package>\tl_greplace_all:Nnn \DocFileDate { / } { - }
%</package|fntef>
%
%<*internal>
\fi
%</internal>
%
%<*driver>
\documentclass[11pt, letterpaper]{l3doc}
\providecommand*\LuaLaTeX{\hologo{LuaLaTeX}}
\providecommand*\XeLaTeX{\hologo{XeLaTeX}}
\usepackage[color = 00498F, ratio = .625]{docext}
\usepackage{fontspec}
\setmainfont{texgyrepagella}[
  Extension      = .otf,
  UprightFont    = *-regular,
  BoldFont       = *-bold,
  ItalicFont     = *-italic,
  BoldItalicFont = *-bolditalic]
\setsansfont{texgyreheros}[
  Extension      = .otf,
  UprightFont    = *-regular,
  BoldFont       = *-bold,
  ItalicFont     = *-italic,
  BoldItalicFont = *-bolditalic]
\setmonofont{cmun}[
  Extension      = .otf,
  UprightFont    = *btl,
  BoldFont       = *tb,
  ItalicFont     = *bto,
  BoldItalicFont = *tx,
  HyphenChar     = None]
\usepackage{unicode-math}
\setmathfont{texgyrepagella-math.otf}
\usepackage{LuaCJK}
\setCJKmainfont{FandolSong-Regular.otf}[
  BoldFont   = FandolSong-Bold.otf,
  ItalicFont = FandolKai-Regular.otf
]
\setCJKsansfont{FandolHei-Regular.otf}[BoldFont = FandolHei-Bold.otf]
\setCJKmonofont{FandolFang-Regular.otf}
\setCJKfamilyfont{zhsong}{FandolSong-Regular.otf}[BoldFont = FandolSong-Bold.otf]
\setCJKfamilyfont{zhhei}{FandolHei-Regular.otf}[BoldFont = FandolHei-Bold.otf]
\setCJKfamilyfont{zhkai}{FandolKai-Regular.otf}
\setCJKfamilyfont{zhfs}{FandolFang-Regular.otf}
\usepackage{LuaCJKfntef}
\LuaCJKsetup{CJKmath=true}
\setCJKmathfont{FandolSong-Regular.otf}
\begin{document}
  \DocInput{\jobname.dtx}
  \clearpage
  \PrintIndex
\end{document}
%</driver>
% \fi
%
% \title{\bfseries The \pkg{LuaCJK} package}
% \author{Mingyu Xia\thanks{\url{https://github.com/myhsia/LuaCJK}}}
% \date{Released \DocFileDate\quad \texttt{\DocFileVersion}}
% \maketitle
%
% \begin{abstract}
%   \pkg{LuaCJK} provides Chinese typesetting for \LuaLaTeX.  It is inspired
%   by \pkg{xeCJK}, but uses \LuaTeX\ node callbacks instead of \XeTeX
%   character classes.  This manual documents the first development release.
% \end{abstract}
%
% \tableofcontents
%
% \begin{documentation}
%
% \section{Introduction}
%
% \pkg{LuaCJK} is a \LuaTeX-only package for Chinese typesetting.  Its public
% interface follows the familiar \pkg{xeCJK} commands:
% \cs{setCJKmainfont}, \cs{setCJKsansfont}, \cs{setCJKmonofont},
% \cs{setCJKfamilyfont}, \cs{CJKfamily}, and \cs{LuaCJKsetup}.
%
% The implementation does not use \XeTeX's inter-character token mechanism.
% Instead, LuaCJK records the current CJK font in a \LuaTeX\ attribute and
% replaces CJK glyph nodes before line breaking and box packing.
%
% \section{Loading}
%
% \begin{quote}
%   \cs{usepackage}\oarg{options}"{LuaCJK}"
% \end{quote}
%
% LuaCJK requires \LuaLaTeX.  It conflicts with \pkg{luatexja}: loading
% either package before or after the other raises an error.
%
% \begin{function}{\LuaCJKversion}
%   Print the version of LuaCJK.
% \end{function}
%
% \section{Font selection}
%
% The main font commands take an optional list of \pkg{fontspec} options.
% Options may be written before or after the font name.
%
% \begin{function}{\setCJKmainfont,\setCJKsansfont,\setCJKmonofont}
%   \begin{syntax}
%     \tn{setCJKmainfont} \oarg{options} \marg{font} \oarg{options}
%   \end{syntax}
%   Set the CJK font for the roman, sans serif, and monospaced families.
% \end{function}
%
% \iffalse
%<*example>
% \fi
\begin{codemo}
\CJKfamily{zhsong}中文宋体\quad
\CJKfamily {zhhei}中文黑体\\
\CJKfamily {zhkai}中文楷体\quad
\CJKfamily  {zhfs}中文仿宋
\end{codemo}
% \iffalse
%</example>
% \fi
%
% \begin{function}{\setCJKfamilyfont}
%   \begin{syntax}
%     \tn{setCJKfamilyfont} \marg{family} \oarg{options} \marg{font} \oarg{options}
%   \end{syntax}
%   Define a named CJK family.
% \end{function}
%
% \begin{function}{\CJKfamily}
%   \begin{syntax}
%     \tn{CJKfamily} \marg{family}
%   \end{syntax}
%   Switch to a named CJK family.
% \end{function}
%
% \begin{function}{\newCJKfontfamily}
%   \begin{syntax}
%     \tn{newCJKfontfamily} \oarg{family} \meta{command} \oarg{options} \marg{font}
%   \end{syntax}
%   Define a command that switches to a newly declared family.  The command is
%   local to the current group, like \cs{newfontfamily}.
% \end{function}
%
% \begin{function}{\CJKfontspec}
%   \begin{syntax}
%     \tn{CJKfontspec} \oarg{options} \marg{font}
%   \end{syntax}
%   Select a CJK font for the current group.
% \end{function}
%
% \begin{function}{\songti, \heiti, \fangsong, \kaishu, \lishu, \youyuan}
%   Shortcuts for the families named \texttt{zhsong}, \texttt{zhhei},
%   \texttt{zhfs}, \texttt{zhkai}, \texttt{zhli}, and \texttt{zhyou}.
% \end{function}
%
% \begin{variable}
%   {\CJKrmdefault, \CJKsfdefault, \CJKttdefault, \CJKfamilydefault}
%   The CJK families selected by \tn{rmfamily}, \tn{sffamily}, and
%   \tn{ttfamily}.
% \end{variable}

% \begin{function}{\defaultCJKfontfeatures}
%   \begin{syntax}
%     \tn{defaultCJKfontfeatures} \marg{font features}
%   \end{syntax}
%   Set the features prepended to every CJK family registered afterwards.
%   The initial value is |Script = CJK|.  Families registered before the
%   call keep the features they were given.
% \end{function}
%
% \begin{function}{\addCJKfontfeatures}
%   \begin{syntax}
%     \tn{addCJKfontfeatures} \oarg{ignored} \marg{font features}
%   \end{syntax}
%   Add features to the current CJK family for the rest of the group, for
%   example |\addCJKfontfeatures{Scale = 1.2}|.  A leading \texttt{*} is
%   accepted for compatibility with \pkg{xeCJK}, where it also covers the
%   sub-blocks; LuaCJK has no blocks, so both forms act on the same font.
% \end{function}
%
% \section{Options}
%
% \begin{function}{\LuaCJKsetup}
%   \begin{syntax}
%     \tn{LuaCJKsetup} \marg{key-value list}
%   \end{syntax}
%   Configure LuaCJK.  The settings are local to the current group.
% \end{function}
%
% The same keys may be given as package options:
% \begin{quote}
%   |\usepackage[PunctStyle = banjiao]{LuaCJK}|
% \end{quote}
%
% \section{Spacing}
%
% \begin{keyval}
%   \item[\keys{CJKglue}]\val{glue}: Glue inserted between adjacent CJK
%   characters.  The initial value is
%   \texttt{0pt plus 0.08\cs{baselineskip}}.
%   \item[\keys{CJKecglue}]\val{glue}: Glue inserted between CJK and Latin
%   characters.  The initial value is a normal interword space.
%   \item[\keys{CJKspace}]\val{\TFF}: Whether to keep source spaces between
%   CJK characters.
%   \item[\keys{xCJKecglue}]\val{\TFF}: Whether to replace source spaces
%   between CJK and Latin characters with \texttt{CJKecglue}.
% \end{keyval}
%
% \iffalse
%<*example>
% \fi
\begin{codemo}
中文 English 中文
\end{codemo}
% \iffalse
%</example>
% \fi
%
% \iffalse
%<*example>
% \fi
\begin{codemo}
\begingroup
  \LuaCJKsetup{CJKspace = true}
  中 文 中 文
\endgroup
\end{codemo}
% \iffalse
%</example>
% \fi
%
% \begin{function}{\normalspacedchars}
%   \begin{syntax}
%     \tn{normalspacedchars} \marg{character list}
%   \end{syntax}
%   The listed characters get no glue on either side, so that a hyphen or a
%   slash inside Latin text does not pull the surrounding text apart.  The
%   initial list is |/|, |\|, and |-|.  The setting is global.
% \end{function}
%
% \section{Suspending the adjustments}
%
% \begin{function}{\LuaCJKdisable, \LuaCJKenable}
%   \begin{syntax}
%     \tn{LuaCJKdisable} \\
%     \tn{LuaCJKenable}
%   \end{syntax}
%   \tn{LuaCJKdisable} stops the typographic adjustments: no \texttt{CJKglue}
%   or \texttt{CJKecglue} is inserted, punctuation is not compressed, no
%   line-breaking prohibitions are added, punctuation is not hung, and CJK
%   characters in math mode are left alone.  CJK characters are still
%   typeset in the CJK font, so text stays readable.  The switch is local to
%   the current group; \tn{LuaCJKenable} turns the adjustments back on.
% \end{function}
%
% This is the hook verbatim-like environments use, and it is also useful when
% measuring raw node lists.
%
% \iffalse
%<*example>
% \fi
\begin{codemo}
\begingroup
  \LuaCJKdisable
  中 文
\endgroup
\end{codemo}
% \iffalse
%</example>
% \fi
%
% \begin{keyval}
%   \item[\keys{LocalConfig}]\val{\TTF \textbar{} \meta{name} \textbar{}
%     \texttt{false}}
%   Read local settings from a configuration file.  \texttt{true} loads
%   \file{luacjk.cfg}, \meta{name} loads \file{luacjk-\meta{name}.cfg}, and
%   \texttt{false} (the initial value) loads nothing.  The file is read at
%   the end of the package, so it may use every command LuaCJK defines.  A
%   missing file only warns.
% \end{keyval}
%
% \begin{keyval}
%   \item[\keys{Verb}]\val{\TTF \textbar{} \texttt{env} \textbar{}
%     \texttt{env+}}
%   What to do inside \tn{verb} and the verbatim-like environments
%   (\env{verbatim}, \env{Verbatim}, \env{lstlisting}, and their
%   relatives).  \texttt{true} suspends the adjustments in both;
%   \texttt{env} keeps them inside the environments but not in \tn{verb},
%   which preserves column alignment when the code contains CJK text;
%   \texttt{env+} keeps them everywhere; \texttt{false} leaves verbatim
%   material entirely alone.  Since \texttt{env+} and \texttt{false} both
%   leave the adjustments on, they behave the same in LuaCJK.
% \end{keyval}
%
% \section{CJK in mathematics}
%
% \begin{keyval}
%   \item[\keys{CJKmath}]\val{\TFF}: Whether CJK characters may be typed
%   directly in math mode.  When it is
%   on, such a character is typeset as a box in the CJK math font instead of
%   being looked up in the math font.  Enable it with
%   \cs{LuaCJKsetup}\marg{CJKmath = true} or with the package option
%   |CJKmath = true|.  Without it a CJK character in math mode is looked up
%   in the math font and usually disappears, so LuaCJK issues a one-time
%   warning in that case.
% \end{keyval}
%
% \begin{function}{\setCJKmathfont}
%   \begin{syntax}
%     \tn{setCJKmathfont} \oarg{options} \marg{font} \oarg{options}
%   \end{syntax}
%   Select the CJK font used for characters in math mode.  Without this
%   command the CJK main font is used.  This manual enables \texttt{CJKmath}
%   and sets this font in its preamble.
% \end{function}
%
% \iffalse
%<*example>
% \fi
\begin{codemo}
$ x + 中文 $ and 中文
\end{codemo}
% \iffalse
%</example>
% \fi
%
% \section{Punctuation}
%
% \begin{keyval}
%   \item[\keys{PunctStyle}]\val{style}: Select a punctuation style.  The
%   implemented styles are
%   \texttt{quanjiao} (the default), \texttt{plain}, \texttt{banjiao},
%   \texttt{kaiming}, \texttt{CCT}, and \texttt{hangmobanjiao}.  An unknown
%   name issues a warning and falls back to \texttt{plain}.
%   \item[\keys{PunctKinsoku}]\val{\TTF}: Whether to apply line-breaking
%   prohibitions for punctuation.  A closing
%   or middle punctuation cannot start a line, an opening punctuation cannot
%   end a line, and a run of long punctuation is kept together.  The
%   prohibitions are independent of \texttt{PunctStyle}, so they stay active
%   with |PunctStyle = plain|.
% \end{keyval}
%
% \begin{function}{\LuaCJKsetwidth}
%   \begin{syntax}
%     \tn{LuaCJKsetwidth} \marg{punctuation list} \marg{length}
%   \end{syntax}
%   Give the listed punctuation marks a fixed advance, overriding the
%   selected \texttt{PunctStyle}.  A leading \texttt{*} is accepted for
%   compatibility with \pkg{xeCJK}, where it sets the width used at the
%   start or end of a line; that case is not implemented and warns once.
% \end{function}
%
% \begin{function}{\LuaCJKsetkern}
%   \begin{syntax}
%     \tn{LuaCJKsetkern} \marg{before} \marg{after} \marg{length}
%   \end{syntax}
%   Set the gap between two adjacent punctuation marks, overriding the
%   style.  Only the first character of each argument is used.
% \end{function}
%
% \begin{function}{\LuaCJKDeclarePunctStyle}
%   \begin{syntax}
%     \tn{LuaCJKDeclarePunctStyle} \marg{name} \marg{key-value list}
%   \end{syntax}
%   Define a punctuation style.  The available keys are
%   \texttt{compress-adjacent}, \texttt{adjacent-ratio}, and
%   \texttt{keep-open-close}.
% \end{function}
%
% \iffalse
%<*example>
% \fi
\begin{codemo}
字，，字
\end{codemo}
% \iffalse
%</example>
% \fi

% \section{Font effects}
%
% The companion package \pkg{LuaCJKfntef} provides underlining, emphasis
% dots, and strike-through.  It builds on \pkg{lua-ul}, so the decorated text
% may contain arbitrary commands, CJK and Latin text, and may still be
% hyphenated and broken between lines.
%
% \begin{function}{\CJKunderline}
%   \begin{syntax}
%     \tn{CJKunderline} \oarg{options} \marg{text}
%   \end{syntax}
%   Underline \meta{text}.  The \meta{options} are passed to the \pkg{lua-ul}
%   underline type, so |top| and |bottom| may be used to tune the position.
%   A leading \texttt{*} or \texttt{-} is accepted for compatibility with
%   \pkg{xeCJKfntef} and ignored.
% \end{function}
%
% \begin{function}{\CJKunderdot}
%   \begin{syntax}
%     \tn{CJKunderdot} \marg{text}
%   \end{syntax}
%   Put emphasis dots under \meta{text}.
% \end{function}
%
% \begin{function}{\CJKsout}
%   \begin{syntax}
%     \tn{CJKsout} \oarg{color} \marg{text}
%   \end{syntax}
%   Strike \meta{text} through.  The optional argument is the color of the
%   rule.
% \end{function}
%
% \iffalse
%<*example>
% \fi
\begin{codemo}
\CJKunderline{中文 test 中文}\\
\CJKunderdot{中文 test 中文}\\
\CJKsout{中文 test 中文}
\end{codemo}
% \iffalse
%</example>
% \fi
%
% \section{Current limitations}
%
% The first release covers the font and spacing core, the
% \texttt{quanjiao} and \texttt{plain} punctuation styles, and the basic
% line-breaking prohibitions for punctuation, hanging punctuation, and direct
% CJK input in math mode, and the font effects of \pkg{LuaCJKfntef}.
% Line-end compression, fallback fonts, alternate fonts, wave and double
% underlines, and a \pkg{ctex} backend are not yet implemented.
%
% \end{documentation}
%
% \StopEventually{\clearpage\PrintIndex}
%
% \begin{implementation}
%
% \section{Implementation}
%
% \begin{macrocode}
%<*package>
%    \end{macrocode}
%
% \begin{macrocode}
%<@@=LuaCJK>
%    \end{macrocode}
%
% \subsection{Engine check and package loading}
%
% LuaCJK only works with \LuaTeX.  The engine test runs before any other
% package is loaded, so that a wrong engine stops with a clear message.
%    \begin{macrocode}
\sys_if_engine_luatex:F
  {
    \msg_new:nnn { LuaCJK } { engine }
      { LuaCJK~requires~LuaTeX. }
    \msg_critical:nn { LuaCJK } { engine }
    \tex_endinput:D
  }
\RequirePackage { l3keys2e }
\RequirePackage { fontspec }
\RequirePackage { xparse }
%    \end{macrocode}
%
% \begin{macro}{\LuaCJKversion}
%   The version of the package.
%   \begin{variable}[internal]{\c_LuaCJK_version_tl}
%     The version string, taken from the package metadata.
%    \begin{macrocode}
\tl_const:Nx \c_LuaCJK_version_tl { \ExplFileVersion }
%    \end{macrocode}
%   \end{variable}
%    \begin{macrocode}
\cs_new:Npn \LuaCJKversion { \tl_use:N \c_LuaCJK_version_tl }
%    \end{macrocode}
% \end{macro}
%
% \subsection{Conflicting packages}
%
% \pkg{luatexja} installs its own CJK font and spacing layer and works on the
% same callbacks as LuaCJK, so the two are rejected whatever the loading
% order.  \pkg{xeCJK} needs no check here: it stops by itself on any engine
% other than \XeTeX.
%    \begin{macrocode}
\msg_new:nnn { LuaCJK } { luatexja-conflict }
  {
    LuaCJK~and~LuaTeX-ja~must~not~be~loaded~together.\\
    Load~only~one~of~them.
  }
\cs_if_exist:cT { ver@luatexja.sty }
  { \msg_critical:nn { LuaCJK } { luatexja-conflict } }
\AddToHook { package/luatexja/before }
  { \msg_critical:nn { LuaCJK } { luatexja-conflict } }
%    \end{macrocode}
%
% \subsection{Attributes and registers}
%
% \begin{variable}[internal]
%   {
%     \g_LuaCJK_font_attr, \c_LuaCJK_font_attr_int,
%     \g_LuaCJK_cjk_space_attr, \c_LuaCJK_cjk_space_attr_int,
%     \g_LuaCJK_xecglue_attr, \c_LuaCJK_xecglue_attr_int,
%     \g_LuaCJK_inserted_attr, \c_LuaCJK_inserted_attr_int,
%     \g_LuaCJK_math_font_attr, \c_LuaCJK_math_font_attr_int,
%     \g_LuaCJK_disabled_attr, \c_LuaCJK_disabled_attr_int
%   }
%   The attributes record, per node, the CJK font, the space and glue
%   policies, whether the node was inserted by LuaCJK, the CJK math font,
%   and whether the adjustments are suspended.  The integer constants hold
%   the allocated attribute numbers for the Lua side.
%    \begin{macrocode}
\newattribute \g_LuaCJK_font_attr
\int_const:Nn \c_LuaCJK_font_attr_int { \the \allocationnumber }
\newattribute \g_LuaCJK_cjk_space_attr
\int_const:Nn \c_LuaCJK_cjk_space_attr_int { \the \allocationnumber }
\newattribute \g_LuaCJK_xecglue_attr
\int_const:Nn \c_LuaCJK_xecglue_attr_int { \the \allocationnumber }
\newattribute \g_LuaCJK_inserted_attr
\int_const:Nn \c_LuaCJK_inserted_attr_int { \the \allocationnumber }
\newattribute \g_LuaCJK_math_font_attr
\int_const:Nn \c_LuaCJK_math_font_attr_int { \the \allocationnumber }
\newattribute \g_LuaCJK_disabled_attr
\int_const:Nn \c_LuaCJK_disabled_attr_int { \the \allocationnumber }
\g_LuaCJK_math_font_attr = 0
\g_LuaCJK_cjk_space_attr = 0
\g_LuaCJK_xecglue_attr = 0
%    \end{macrocode}
% \end{variable}
%
% \begin{variable}[internal]{\l_LuaCJK_ccglue_skip, \l_LuaCJK_ecglue_skip}
%   The glue used between CJK characters and between CJK and Latin text.
%    \begin{macrocode}
\skip_new:N \l_LuaCJK_ccglue_skip
\skip_new:N \l_LuaCJK_ecglue_skip
\skip_set:Nn \l_LuaCJK_ccglue_skip { 0pt plus 0.08 \tex_baselineskip:D }
\skip_set:Nn \l_LuaCJK_ecglue_skip
  { 0.3333em plus 0.1667em minus 0.1111em }
%    \end{macrocode}
% \end{variable}
%
% \subsection{Punctuation styles}
%
% \begin{variable}[internal]{\g_LuaCJK_punct_style_prop}
%   The names of the styles LuaCJK knows about, including those defined by
%   \cs{LuaCJKDeclarePunctStyle}.
%    \begin{macrocode}
\prop_new:N \g_LuaCJK_punct_style_prop
\prop_gput:Nnn \g_LuaCJK_punct_style_prop { quanjiao } { }
\prop_gput:Nnn \g_LuaCJK_punct_style_prop { plain } { }
\prop_gput:Nnn \g_LuaCJK_punct_style_prop { banjiao } { }
\prop_gput:Nnn \g_LuaCJK_punct_style_prop { kaiming } { }
\prop_gput:Nnn \g_LuaCJK_punct_style_prop { CCT } { }
\prop_gput:Nnn \g_LuaCJK_punct_style_prop { hangmobanjiao } { }
%    \end{macrocode}
% \end{variable}
%
% \begin{macro}[internal]{\LuaCJK_set_punct_style:n}
%   Store the selected style, or warn and fall back to \texttt{plain} when
%   the name is unknown.
%   \begin{variable}[internal]{\l_LuaCJK_punct_style_tl}
%     The style name read by the Lua side.
%    \begin{macrocode}
\tl_new:N \l_LuaCJK_punct_style_tl
\tl_set:Nn \l_LuaCJK_punct_style_tl { quanjiao }
%    \end{macrocode}
%   \end{variable}
%    \begin{macrocode}
\msg_new:nnn { LuaCJK } { punct-style-unknown }
  {
    PunctStyle~`#1'~is~not~implemented.\\
    Using~`plain'~instead.
  }
\cs_new_protected:Npn \LuaCJK_set_punct_style:n #1
  {
    \prop_if_in:NnTF \g_LuaCJK_punct_style_prop {#1}
      { \tl_set:Nn \l_LuaCJK_punct_style_tl {#1} }
      {
        \msg_warning:nnn { LuaCJK } { punct-style-unknown } {#1}
        \tl_set:Nn \l_LuaCJK_punct_style_tl { plain }
      }
  }
%    \end{macrocode}
% \end{macro}
%
% \begin{macro}{\LuaCJKsetwidth}
%   Give specific punctuation marks a fixed advance, overriding the
%   selected style.  The star form of \pkg{xeCJK}, which sets the width used
%   at the start or end of a line, has no counterpart yet and only warns.
%   \begin{variable}[internal]{\l_LuaCJK_bound_width_warned_bool}
%     Guards the one-time warning for the star form.
%    \begin{macrocode}
\bool_new:N \l_LuaCJK_bound_width_warned_bool
\msg_new:nnn { LuaCJK } { bound-width-unsupported }
  {
    The~star~form~of~\token_to_str:N \LuaCJKsetwidth,~which~sets~the~width~
    at~the~start~or~end~of~a~line,~is~not~implemented.\\
    The~width~is~applied~to~every~position~instead.
  }
%    \end{macrocode}
%   \end{variable}
%    \begin{macrocode}
\NewDocumentCommand \LuaCJKsetwidth { s m m }
  {
    \IfBooleanT {#1}
      {
        \bool_if:NF \l_LuaCJK_bound_width_warned_bool
          {
            \bool_set_true:N \l_LuaCJK_bound_width_warned_bool
            \msg_warning:nn { LuaCJK } { bound-width-unsupported }
          }
      }
    \lua_now:e
      {
        LuaCJK.set_punct_width
          ( "\lua_escape:e {#2}" , \dim_to_decimal_in_sp:n {#3} )
      }
  }
%    \end{macrocode}
% \end{macro}
%
% \begin{macro}{\LuaCJKsetkern}
%   Set the gap between two adjacent punctuation marks, overriding the
%   selected style.  Only the first character of each argument is used.
%    \begin{macrocode}
\NewDocumentCommand \LuaCJKsetkern { m m m }
  {
    \lua_now:e
      {
        LuaCJK.set_punct_kern
          ( "\lua_escape:e {#1}" , "\lua_escape:e {#2}"
          , \dim_to_decimal_in_sp:n {#3} )
      }
  }
%    \end{macrocode}
% \end{macro}
%
% \begin{macro}{\LuaCJKDeclarePunctStyle}
%   Define a style from a small set of keys and hand the result to Lua.
%   \begin{variable}[internal]
%     {\l_LuaCJK_punct_compress_bool, \l_LuaCJK_punct_ratio_fp,
%      \l_LuaCJK_punct_keep_bool}
%     The values collected from the keys, with the defaults used when a key
%     is not given.
%    \begin{macrocode}
\bool_new:N \l_LuaCJK_punct_compress_bool
\fp_new:N \l_LuaCJK_punct_ratio_fp
\bool_new:N \l_LuaCJK_punct_keep_bool
%    \end{macrocode}
%   \end{variable}
%    \begin{macrocode}
\keys_define:nn { LuaCJK / punct-style }
  {
    compress-adjacent .bool_set:N = \l_LuaCJK_punct_compress_bool,
    adjacent-ratio .fp_set:N = \l_LuaCJK_punct_ratio_fp,
    keep-open-close .bool_set:N = \l_LuaCJK_punct_keep_bool,
  }
\NewDocumentCommand \LuaCJKDeclarePunctStyle { m m }
  {
    \bool_set_true:N \l_LuaCJK_punct_compress_bool
    \fp_set:Nn \l_LuaCJK_punct_ratio_fp { 0.5 }
    \bool_set_true:N \l_LuaCJK_punct_keep_bool
    \keys_set:nn { LuaCJK / punct-style } {#2}
    \prop_gput:Nnn \g_LuaCJK_punct_style_prop {#1} { }
    \lua_now:e
      {
        LuaCJK.define_punct_style(
          "\lua_escape:e {#1}",
          \bool_if:NTF \l_LuaCJK_punct_compress_bool { true } { false },
          \fp_use:N \l_LuaCJK_punct_ratio_fp,
          \bool_if:NTF \l_LuaCJK_punct_keep_bool { true } { false } )
      }
  }
%    \end{macrocode}
% \end{macro}
%
% \subsection{Options}
%
% \begin{macro}{\LuaCJKsetup}
%   The same key path serves \cs{LuaCJKsetup} and the package options.
%   \begin{variable}[internal]
%     {\l_LuaCJK_punct_kinsoku_tl, \l_LuaCJK_cjk_math_tl, \l_LuaCJK_verb_mode_tl}
%     The switches read by the Lua side, plus the \texttt{Verb} setting.
%    \begin{macrocode}
\tl_new:N \l_LuaCJK_punct_kinsoku_tl
\tl_set:Nn \l_LuaCJK_punct_kinsoku_tl { 1 }
\tl_new:N \l_LuaCJK_cjk_math_tl
\tl_set:Nn \l_LuaCJK_cjk_math_tl { 0 }
\tl_new:N \l_LuaCJK_verb_mode_tl
\tl_set:Nn \l_LuaCJK_verb_mode_tl { true }
\tl_new:N \l_LuaCJK_local_config_tl
%    \end{macrocode}
%   \end{variable}
%    \begin{macrocode}
\keys_define:nn { LuaCJK / options }
  {
    CJKglue .code:n =
      {
        \skip_set:Nn \l_LuaCJK_ccglue_skip {#1}
      },
    CJKecglue .code:n =
      {
        \skip_set:Nn \l_LuaCJK_ecglue_skip {#1}
      },
    CJKspace .choice:,
    CJKspace / true  .code:n = { \g_LuaCJK_cjk_space_attr = 1 },
    CJKspace / false .code:n = { \g_LuaCJK_cjk_space_attr = 0 },
    CJKspace .default:n = { true },
    CJKspace .initial:n = { false },
    xCJKecglue .choice:,
    xCJKecglue / true  .code:n = { \g_LuaCJK_xecglue_attr = 1 },
    xCJKecglue / false .code:n = { \g_LuaCJK_xecglue_attr = 0 },
    xCJKecglue .default:n = { true },
    xCJKecglue .initial:n = { false },
    PunctStyle .code:n = { \LuaCJK_set_punct_style:n {#1} },
    PunctStyle .default:n = { quanjiao },
    PunctKinsoku .choice:,
    PunctKinsoku / true  .code:n = { \tl_set:Nn \l_LuaCJK_punct_kinsoku_tl { 1 } },
    PunctKinsoku / false .code:n = { \tl_set:Nn \l_LuaCJK_punct_kinsoku_tl { 0 } },
    PunctKinsoku .default:n = { true },
    CJKmath .choice:,
    CJKmath / true  .code:n = { \tl_set:Nn \l_LuaCJK_cjk_math_tl { 1 } },
    CJKmath / false .code:n = { \tl_set:Nn \l_LuaCJK_cjk_math_tl { 0 } },
    CJKmath .default:n = { true },
    Verb .choice:,
    Verb / true  .code:n = { \tl_set:Nn \l_LuaCJK_verb_mode_tl { true } },
    Verb / env   .code:n = { \tl_set:Nn \l_LuaCJK_verb_mode_tl { env } },
    Verb / env+  .code:n = { \tl_set:Nn \l_LuaCJK_verb_mode_tl { env+ } },
    Verb / false .code:n = { \tl_set:Nn \l_LuaCJK_verb_mode_tl { false } },
    Verb .default:n = { true },
    LocalConfig .code:n =
      {
        \str_case:nnF {#1}
          {
            { true }  { \tl_set:Nn \l_LuaCJK_local_config_tl { luacjk.cfg } }
            { false } { \tl_clear:N \l_LuaCJK_local_config_tl }
          }
          { \tl_set:Nn \l_LuaCJK_local_config_tl { luacjk-#1.cfg } }
      },
    LocalConfig .default:n = { true },
    unknown .code:n =
      {
        \msg_warning:nne { LuaCJK } { unknown-option }
          { \l_keys_key_str }
      },
  }
\msg_new:nnn { LuaCJK } { unknown-option }
  { Unknown~option~`#1'~is~being~ignored. }
\NewDocumentCommand \LuaCJKsetup { m }
  { \keys_set:nn { LuaCJK / options } {#1} }
\cs_if_exist:NTF \ProcessKeyOptions
  { \ProcessKeyOptions [ LuaCJK / options ] }
  { \ProcessKeysOptions { LuaCJK / options } }
%    \end{macrocode}
% \end{macro}
%
% \subsection{Suspending the adjustments}
%
% \begin{macro}{\LuaCJKdisable, \LuaCJKenable}
%   Set the node attribute that makes the Lua callbacks skip the
%   adjustments.  The attribute follows the \TeX{} grouping.
%    \begin{macrocode}
\NewDocumentCommand \LuaCJKdisable { }
  { \g_LuaCJK_disabled_attr = 1 }
\NewDocumentCommand \LuaCJKenable { }
  { \g_LuaCJK_disabled_attr = 0 }
%    \end{macrocode}
% \end{macro}
%
% \begin{macro}[internal]{\LuaCJK_verbatim@font_hook:}
%   Verbatim-like material must usually be left alone.  \tn{verb} opens its
%   own group and selects \tn{verbatim@font} inside it, so extending that
%   command covers \tn{verb} on its own.  The hook consults the current
%   \texttt{Verb} setting, so that \cs{LuaCJKsetup} can change the behaviour
%   after the package is loaded.
%    \begin{macrocode}
\cs_new_protected:Npn \LuaCJK_verbatim@font_hook:
  {
    \tl_if_eq:VnTF \l_LuaCJK_verb_mode_tl { true }
      { \LuaCJKdisable }
      {
        \tl_if_eq:VnT \l_LuaCJK_verb_mode_tl { env }
          { \LuaCJKdisable }
      }
  }
\cs_if_exist:NT \verbatim@font
  {
    \cs_gset_protected:Npx \verbatim@font
      {
        \exp_not:o { \verbatim@font }
        \exp_not:N \LuaCJK_verbatim@font_hook:
      }
  }
%    \end{macrocode}
% \end{macro}
%
% \begin{macro}[internal]{\LuaCJK_verbatim_env_begin:, \LuaCJK_verbatim_env_end:}
%   Environments that do not go through \tn{verbatim@font} need their own
%   hooks.  Only the \texttt{true} setting suspends the adjustments here:
%   \texttt{env} and \texttt{env+} keep them, which is what keeps code
%   columns aligned when CJK text is used inside the code.
%    \begin{macrocode}
\cs_new_protected:Npn \LuaCJK_verbatim_env_begin:
  {
    \tl_if_eq:VnT \l_LuaCJK_verb_mode_tl { true }
      { \LuaCJKdisable }
  }
\cs_new_protected:Npn \LuaCJK_verbatim_env_end:
  {
    \tl_if_eq:VnT \l_LuaCJK_verb_mode_tl { true }
      { \LuaCJKenable }
  }
\clist_map_inline:nn
  {
    verbatim , verbatim* , Verbatim , Verbatim* ,
    BVerbatim , LVerbatim , lstlisting
  }
  {
    \AddToHook { env / #1 / begin } { \LuaCJK_verbatim_env_begin: }
    \AddToHook { env / #1 / end }   { \LuaCJK_verbatim_env_end: }
  }
%    \end{macrocode}
% \end{macro}
%
% \subsection{CJK families}
%
% \begin{macro}[internal]{\LuaCJK_set_family:nnn}
%   Define a family through \pkg{fontspec} and remember the resulting
%   family name.  The family is only registered; the switch happens in
%   \cs{LuaCJK_switch_family:n}.
%   \begin{variable}[internal]{\g_LuaCJK_family_prop, \l_LuaCJK_family_tl}
%     The registered families, and the family name currently selected.
%    \begin{macrocode}
\prop_new:N \g_LuaCJK_family_prop
\tl_new:N \l_LuaCJK_family_tl
\tl_set:Nn \l_LuaCJK_family_tl { zhsong }
%    \end{macrocode}
%   \end{variable}
%   \begin{variable}[internal]{\l_LuaCJK_fontspec_tl, \l_LuaCJK_family_name_tl, \l_LuaCJK_tmpa_tl}
%     Scratch variables holding the generated \pkg{fontspec} family name,
%     the family name looked up in the property list, and a temporary
%     value.
%    \begin{macrocode}
\tl_new:N \l_LuaCJK_fontspec_tl
\tl_new:N \l_LuaCJK_family_name_tl
\tl_new:N \l_LuaCJK_tmpa_tl
\tl_new:N \l_LuaCJK_registered_name_tl
%    \end{macrocode}
%   \end{variable}
%   \begin{variable}[internal]
%     {\g_LuaCJK_font_name_prop, \g_LuaCJK_font_feature_prop,
%      \l_LuaCJK_default_feature_tl}
%     The font name and the explicit options of each family, kept so that
%     \cs{addCJKfontfeatures} can register the family again, and the
%     features prepended to every family.
%    \begin{macrocode}
\prop_new:N \g_LuaCJK_font_name_prop
\prop_new:N \g_LuaCJK_font_feature_prop
\tl_new:N \l_LuaCJK_default_feature_tl
\tl_set:Nn \l_LuaCJK_default_feature_tl { Script = CJK }
\tl_new:N \l_LuaCJK_feature_tl
\cs_generate_variant:Nn \fontspec_gset_family:Nnn { NVV }
%    \end{macrocode}
%   \end{variable}
%    \begin{macrocode}
\cs_new_protected:Npn \LuaCJK_set_family:nnn #1#2#3
  {
    \tl_set:Ne \l_LuaCJK_feature_tl
      {
        \exp_not:V \l_LuaCJK_default_feature_tl
        \tl_if_blank:nF {#2} { , \exp_not:n {#2} }
      }
    \prop_gput:NnV \g_LuaCJK_font_feature_prop {#1} \l_LuaCJK_feature_tl
    \prop_gput:Nnn \g_LuaCJK_font_name_prop {#1} {#3}
    \LuaCJK_register_family:n {#1}
  }
%    \end{macrocode}
%   \begin{macro}[internal]{\LuaCJK_register_family:n}
%     Build the \pkg{fontspec} family from the stored font name and
%     features, and re-select it if it is the current one.
%    \begin{macrocode}
\cs_new_protected:Npn \LuaCJK_register_family:n #1
  {
    \prop_get:NnNF \g_LuaCJK_font_name_prop {#1} \l_LuaCJK_registered_name_tl
      { \tl_clear:N \l_LuaCJK_registered_name_tl }
    \prop_get:NnNF \g_LuaCJK_font_feature_prop {#1} \l_LuaCJK_feature_tl
      { \tl_clear:N \l_LuaCJK_feature_tl }
    \group_begin:
      \fontspec_gset_family:NVV \l_LuaCJK_fontspec_tl
        \l_LuaCJK_feature_tl \l_LuaCJK_registered_name_tl
      \prop_gput:NnV \g_LuaCJK_family_prop {#1} \l_LuaCJK_fontspec_tl
    \group_end:
    \prop_if_in:NVT \g_LuaCJK_family_prop \l_LuaCJK_family_tl
      { \selectfont }
  }
%    \end{macrocode}
%   \end{macro}
% \end{macro}
%
% \begin{macro}[internal]{\LuaCJK_switch_family:n}
%   Select a registered family and warn once about an unknown name.
%    \begin{macrocode}
\cs_new_protected:Npn \LuaCJK_switch_family:n #1
  {
    \prop_if_in:NnTF \g_LuaCJK_family_prop {#1}
      {
        \tl_set:Nn \l_LuaCJK_family_tl {#1}
        \selectfont
      }
      {
        \seq_if_in:NnF \g_LuaCJK_unknown_family_seq {#1}
          {
            \seq_gput_right:Nn \g_LuaCJK_unknown_family_seq {#1}
            \msg_warning:nnn { LuaCJK } { family-unknown } {#1}
          }
      }
  }
\seq_new:N \g_LuaCJK_unknown_family_seq
\msg_new:nnn { LuaCJK } { family-unknown }
  { Unknown~CJK~family~`#1'~is~being~ignored. }
%    \end{macrocode}
% \end{macro}
%
% \begin{macro}[internal]{\LuaCJK_update_font:}
%   On every \tn{selectfont}, resolve the current family and the math
%   family to font identifiers and store them in the attributes.
%   \begin{variable}[internal]{\l_LuaCJK_selectfont_bool, \g_LuaCJK_fontid_int}
%     A guard against re-entering the update while \tn{selectfont} runs,
%     and the font identifier of the current CJK family.
%    \begin{macrocode}
\bool_new:N \l_LuaCJK_selectfont_bool
\int_new:N \g_LuaCJK_fontid_int
%    \end{macrocode}
%   \end{variable}
%   \begin{variable}[internal]{\l_LuaCJK_math_family_tl, \g_LuaCJK_math_fontid_int}
%     The math family registered by \cs{setCJKmathfont}, and the font
%     identifier resolved from it.
%    \begin{macrocode}
\tl_new:N \l_LuaCJK_math_family_tl
\int_new:N \g_LuaCJK_math_fontid_int
%    \end{macrocode}
%   \end{variable}
%    \begin{macrocode}
\cs_new_protected:Npn \LuaCJK_update_font:
  {
    \bool_if:NF \l_LuaCJK_selectfont_bool
      {
        \bool_set_true:N \l_LuaCJK_selectfont_bool
        \prop_get:NVNTF \g_LuaCJK_family_prop \l_LuaCJK_family_tl \l_LuaCJK_family_name_tl
          {
            \group_begin:
              \exp_args:NV \fontfamily \l_LuaCJK_family_name_tl
              \selectfont
              \int_gset:Nn \g_LuaCJK_fontid_int
                { \tex_fontid:D \tex_font:D }
            \group_end:
            \g_LuaCJK_font_attr = \g_LuaCJK_fontid_int
          }
          { \g_LuaCJK_font_attr = \c_zero_int }
        \tl_if_empty:NTF \l_LuaCJK_math_family_tl
          { \int_gset:Nn \g_LuaCJK_math_fontid_int { \int_use:N \g_LuaCJK_fontid_int } }
          {
            \prop_get:NVNTF \g_LuaCJK_family_prop \l_LuaCJK_math_family_tl
              \l_LuaCJK_tmpa_tl
              {
                \group_begin:
                  \exp_args:NV \fontfamily \l_LuaCJK_tmpa_tl
                  \selectfont
                  \int_gset:Nn \g_LuaCJK_math_fontid_int
                    { \tex_fontid:D \tex_font:D }
                \group_end:
              }
              { \int_gzero:N \g_LuaCJK_math_fontid_int }
          }
        \g_LuaCJK_math_font_attr = \g_LuaCJK_math_fontid_int
        \bool_set_false:N \l_LuaCJK_selectfont_bool
      }
  }
\AddToHook { selectfont } { \LuaCJK_update_font: }
%    \end{macrocode}
% \end{macro}
%
% \subsection{Font selection commands}
%
% \begin{variable}
%   {\CJKrmdefault, \CJKsfdefault, \CJKttdefault, \CJKfamilydefault}
%   The CJK families used by the standard font commands.  The family
%   hooks below re-select the CJK family whenever \tn{rmfamily} and its
%   siblings run.
%    \begin{macrocode}
\tl_new:N \CJKrmdefault
\tl_new:N \CJKsfdefault
\tl_new:N \CJKttdefault
\tl_new:N \CJKfamilydefault
\tl_set:Nn \CJKrmdefault { zhsong }
\tl_set:Nn \CJKsfdefault { zhhei }
\tl_set:Nn \CJKttdefault { zhfs }
\tl_set:Nn \CJKfamilydefault { zhsong }
%    \end{macrocode}
% \end{variable}
%
% \begin{macro}{\defaultCJKfontfeatures}
%   Set the features prepended to every CJK family, replacing the initial
%   \texttt{Script = CJK}.  Families registered earlier keep the value they
%   were given; call \cs{setCJKmainfont} again to apply a new default.
%    \begin{macrocode}
\NewDocumentCommand \defaultCJKfontfeatures { m }
  { \tl_set:Nn \l_LuaCJK_default_feature_tl {#1} }
%    \end{macrocode}
% \end{macro}
%
% \begin{macro}{\addCJKfontfeatures}
%   Add features to the current family for the rest of the group.  The star
%   form is accepted for compatibility; LuaCJK has no concept of blocks, so
%   both forms act on the same font.
%    \begin{macrocode}
\NewDocumentCommand \addCJKfontfeatures { s o m }
  {
    \exp_args:NV \LuaCJK_add_family_features:n
      \l_LuaCJK_family_tl {#3}
  }
%    \end{macrocode}
%   \begin{macro}[internal]{\LuaCJK_add_family_features:n}
%     Append features to the stored ones of a family and register it again.
%    \begin{macrocode}
\cs_new_protected:Npn \LuaCJK_add_family_features:n #1#2
  {
    \prop_get:NnNTF \g_LuaCJK_font_feature_prop {#1} \l_LuaCJK_feature_tl
      {
        \tl_set:Ne \l_LuaCJK_feature_tl
          {
            \exp_not:V \l_LuaCJK_feature_tl
            \tl_if_blank:nF {#2} { , \exp_not:n {#2} }
          }
        \prop_get:NnNF \g_LuaCJK_font_name_prop {#1} \l_LuaCJK_registered_name_tl
          { \tl_clear:N \l_LuaCJK_registered_name_tl }
        \group_begin:
          \fontspec_gset_family:NVV \l_LuaCJK_fontspec_tl
            \l_LuaCJK_feature_tl \l_LuaCJK_registered_name_tl
        \group_end:
        \prop_put:NnV \g_LuaCJK_family_prop {#1} \l_LuaCJK_fontspec_tl
        \prop_put:NnV \g_LuaCJK_font_feature_prop {#1}
          \l_LuaCJK_feature_tl
        \selectfont
      }
      { \msg_warning:nnn { LuaCJK } { family-unknown } {#1} }
  }
%    \end{macrocode}
%   \end{macro}
% \end{macro}
%
% \begin{macro}
%   {\setCJKmainfont, \setCJKsansfont, \setCJKmonofont, \setCJKfamilyfont}
%   The public font commands, all built on
%   \cs{LuaCJK_set_family_with_options:nnnn}.
%    \begin{macrocode}
\NewDocumentCommand \setCJKmainfont { o m o }
  {
    \exp_args:NV \LuaCJK_set_family_with_options:nnnn
      \CJKrmdefault {#1} {#2} {#3}
  }
\NewDocumentCommand \setCJKsansfont { o m o }
  {
    \exp_args:NV \LuaCJK_set_family_with_options:nnnn
      \CJKsfdefault {#1} {#2} {#3}
  }
\NewDocumentCommand \setCJKmonofont { o m o }
  {
    \exp_args:NV \LuaCJK_set_family_with_options:nnnn
      \CJKttdefault {#1} {#2} {#3}
  }
\NewDocumentCommand \setCJKfamilyfont { m o m o }
  { \LuaCJK_set_family_with_options:nnnn {#1} {#2} {#3} {#4} }
%    \end{macrocode}
%   \begin{macro}[internal]{\LuaCJK_set_family_with_options:nnnn}
%     Accept the options before or after the font name.
%    \begin{macrocode}
\cs_new_protected:Npn \LuaCJK_set_family_with_options:nnnn #1#2#3#4
  {
    \tl_if_novalue:nTF {#2}
      {
        \tl_if_novalue:nTF {#4}
          { \LuaCJK_set_family:nnn {#1} { } {#3} }
          { \LuaCJK_set_family:nnn {#1} {#4} {#3} }
      }
      { \LuaCJK_set_family:nnn {#1} {#2} {#3} }
  }
%    \end{macrocode}
%   \end{macro}
% \end{macro}
%
% \begin{macro}{\CJKfamily}
%   Switch to a registered family and gobble the following spaces.
%    \begin{macrocode}
\NewDocumentCommand \CJKfamily { m }
  { \LuaCJK_switch_family:n {#1} \tex_ignorespaces:D }
%    \end{macrocode}
% \end{macro}
%
% \begin{macro}{\newCJKfontfamily}
%   Define a control sequence that switches to a new family, and register
%   the family under the same name.
%    \begin{macrocode}
\NewDocumentCommand \newCJKfontfamily { o m o m }
  {
    \tl_set:Ne \l_LuaCJK_tmpa_tl
      {
        \IfNoValueTF {#1}
          { \cs_to_str:N #2 }
          {#1}
      }
    \cs_set_protected:Npx #2
      { \exp_not:N \LuaCJK_switch_family:n { \exp_not:V \l_LuaCJK_tmpa_tl } }
    \exp_args:NV \LuaCJK_set_family_with_options:nnnn
      \l_LuaCJK_tmpa_tl {#3} {#4} { \NoValue }
  }
%    \end{macrocode}
% \end{macro}
%
% \begin{macro}{\setCJKmathfont}
%   Register a family used for CJK characters in math mode and refresh
%   the attributes at once.
%    \begin{macrocode}
\NewDocumentCommand \setCJKmathfont { o m o }
  {
    \LuaCJK_set_family_with_options:nnnn { LuaCJK/math } {#1} {#2} {#3}
    \tl_set:Nn \l_LuaCJK_math_family_tl { LuaCJK/math }
    \LuaCJK_update_font:
  }
%    \end{macrocode}
% \end{macro}
%
% \begin{macro}{\CJKfontspec}
%   Register a one-off family and switch to it immediately.
%   \begin{variable}[internal]{\g_LuaCJK_fontspec_int}
%     Gives every call its own family name.
%    \begin{macrocode}
\int_new:N \g_LuaCJK_fontspec_int
%    \end{macrocode}
%   \end{variable}
%    \begin{macrocode}
\NewDocumentCommand \CJKfontspec { o m }
  {
    \int_gincr:N \g_LuaCJK_fontspec_int
    \tl_set:Ne \l_LuaCJK_tmpa_tl
      { LuaCJK/fontspec/ \int_use:N \g_LuaCJK_fontspec_int }
    \exp_args:NV \LuaCJK_set_family_with_options:nnnn
      \l_LuaCJK_tmpa_tl {#1} {#2} { \NoValue }
    \exp_args:NV \LuaCJK_switch_family:n \l_LuaCJK_tmpa_tl
    \tex_ignorespaces:D
  }
%    \end{macrocode}
% \end{macro}
%
% \begin{macro}{\songti, \heiti, \fangsong, \kaishu, \lishu, \youyuan}
%   Shortcuts for the usual families, plus the family hooks of the standard
%   font commands.
%    \begin{macrocode}
\cs_new_protected:Npn \songti { \CJKfamily { zhsong } }
\cs_new_protected:Npn \heiti { \CJKfamily { zhhei  } }
\cs_new_protected:Npn \fangsong { \CJKfamily { zhfs   } }
\cs_new_protected:Npn \kaishu { \CJKfamily { zhkai  } }
\cs_new_protected:Npn \lishu { \CJKfamily { zhli   } }
\cs_new_protected:Npn \youyuan { \CJKfamily { zhyou  } }
\AddToHook { cmd/rmfamily/after }
  { \exp_args:NV \LuaCJK_switch_family:n \CJKrmdefault }
\AddToHook { cmd/sffamily/after }
  { \exp_args:NV \LuaCJK_switch_family:n \CJKsfdefault }
\AddToHook { cmd/ttfamily/after }
  { \exp_args:NV \LuaCJK_switch_family:n \CJKttdefault }
%    \end{macrocode}
% \end{macro}
%
% \subsection{Characters without CJK/Latin glue}
%
% \begin{macro}{\normalspacedchars}
%   The characters in the list get no glue on either side, which is what one
%   wants for slashes, backslashes, and hyphens inside Latin text.  The list
%   is kept on the Lua side as a set of code points, so the setting is
%   global: a group does not restore the previous list.
%   \begin{variable}[internal]{\g_LuaCJK_normal_spaced_chars_tl}
%     The list, initially \texttt{/}, \texttt{\c_backslash_str}, and
%     \texttt{-}.
%    \begin{macrocode}
\tl_new:N \g_LuaCJK_normal_spaced_chars_tl
\tl_gset:Nn \g_LuaCJK_normal_spaced_chars_tl
  { / , \c_backslash_str , - }
%    \end{macrocode}
%   \end{variable}
%    \begin{macrocode}
\cs_new_protected:Npn \LuaCJK_set_normal_spaced_chars:
  {
    \lua_now:e
      {
        LuaCJK.set_normal_spaced_chars
          ( "\lua_escape:e { \g_LuaCJK_normal_spaced_chars_tl } " )
      }
  }
\NewDocumentCommand \normalspacedchars { m }
  {
    \tl_gset:Nn \g_LuaCJK_normal_spaced_chars_tl {#1}
    \LuaCJK_set_normal_spaced_chars:
  }
%    \end{macrocode}
% \end{macro}
%
% \subsection{Local configuration}
%
% \begin{macro}{\LuaCJK_load_local_config:}
%   Load the file chosen by the \texttt{LocalConfig} option.  The key is
%   read while the package is being loaded, but the file is only input at
%   the end, so that it can use everything LuaCJK defines.
%    \begin{macrocode}
\cs_new_protected:Npn \LuaCJK_load_local_config:
  {
    \tl_if_empty:NF \l_LuaCJK_local_config_tl
      {
        \exp_args:NV \LuaCJK_load_config_file:n \l_LuaCJK_local_config_tl
      }
  }
\cs_new_protected:Npn \LuaCJK_load_config_file:n #1
  {
    \file_if_exist:nTF {#1}
      { \file_input:n {#1} }
      { \msg_warning:nnn { LuaCJK } { config-missing } {#1} }
  }
\msg_new:nnn { LuaCJK } { config-missing }
  { The~configuration~file~`#1'~was~not~found. }
%    \end{macrocode}
% \end{macro}
%
% \subsection{Loading the Lua core}
%
% The Lua module reads the attribute numbers and installs its callbacks.
% The current font is resolved once the document starts.
%    \begin{macrocode}
\lua_load_module:n { LuaCJK }
\lua_now:e
  {
    LuaCJK.attr_font = \int_use:N \c_LuaCJK_font_attr_int ;
    LuaCJK.attr_cjk_space = \int_use:N \c_LuaCJK_cjk_space_attr_int ;
    LuaCJK.attr_xecglue = \int_use:N \c_LuaCJK_xecglue_attr_int ;
    LuaCJK.attr_inserted = \int_use:N \c_LuaCJK_inserted_attr_int ;
    LuaCJK.attr_math_font = \int_use:N \c_LuaCJK_math_font_attr_int ;
    LuaCJK.attr_disabled = \int_use:N \c_LuaCJK_disabled_attr_int
  }
\LuaCJK_set_normal_spaced_chars:
\AtBeginDocument
  {
    \LuaCJK_update_font:
    \g_LuaCJK_font_attr = \g_LuaCJK_fontid_int
  }
\LuaCJK_load_local_config:
%    \end{macrocode}
%
%    \begin{macrocode}
%</package>
%    \end{macrocode}
%
% \subsection{The Lua entry point}
%
% The package ships two Lua modules.  \file{LuaCJK.lua} only loads the core
% module; all node processing lives in \file{LuaCJK-core.lua}.
% \begin{variable}[internal]{LuaCJK}
%   The table shared by the two modules.  The core module adds its functions
%   and constants to it.
%    \begin{macrocode}
%<*lua>
-- LuaCJK.lua
-- Copyright (C) 2026 by Mingyu Xia <myhsia@outlook.com>
-- This work may be distributed and/or modified under the conditions of the
-- LaTeX Project Public License, version 1.3c or later.
-- Entry point: load the core module.

LuaCJK = LuaCJK or {}

require("LuaCJK-core")
%</lua>
%    \end{macrocode}
% \end{variable}
%
% \subsection{Module state and node types}
%
% The core module starts by creating the shared table and caching the node
% identifiers.  The callbacks run on every list, so these lookups happen
% once.
%    \begin{macrocode}
%<*lua-core>
-- LuaCJK-core.lua
-- Copyright (C) 2026 by Mingyu Xia <myhsia@outlook.com>
-- This work may be distributed and/or modified under the conditions of the
-- LaTeX Project Public License, version 1.3c or later.
-- Core node processing.

LuaCJK = LuaCJK or {}

%    \end{macrocode}
%
% \begin{variable}[internal]
%   {M, GLYPH, GLUE, HLIST, VLIST, KERN, PENALTY, MATH, DISC, WHATSIT,
%    LOCAL_PAR, SPACESKIP, new_node, free_node, has_attribute, set_attribute}
%   Short names for the module table, the node types, and the node
%   functions.
%    \begin{macrocode}
local M = LuaCJK
local GLYPH = node.id("glyph")
local GLUE = node.id("glue")
local HLIST = node.id("hlist")
local VLIST = node.id("vlist")
local KERN = node.id("kern")
local PENALTY = node.id("penalty")
local MATH = node.id("math")
local DISC = node.id("disc")
local WHATSIT = node.id("whatsit")
local LOCAL_PAR = node.id("local_par")
local SPACESKIP = node.subtype("spaceskip") or 13
local new_node = node.new
local free_node = node.free
local has_attribute = node.has_attribute
local set_attribute = node.set_attribute

%    \end{macrocode}
% \end{variable}
%
% \begin{variable}[internal]
%   {LuaCJK.attr_font, LuaCJK.attr_cjk_space, LuaCJK.attr_xecglue,
%    LuaCJK.attr_inserted, LuaCJK.attr_math_font}
%   The attribute numbers.  The \TeX{} side fills them in after loading the
%   module.
%    \begin{macrocode}
M.attr_font = M.attr_font or nil
M.attr_cjk_space = M.attr_cjk_space or nil
M.attr_xecglue = M.attr_xecglue or nil
M.attr_inserted = M.attr_inserted or nil
M.attr_math_font = M.attr_math_font or nil

%    \end{macrocode}
% \end{variable}
%
% \subsection{Character classes}
%
% \begin{macro}[internal]{LuaCJK.is_cjk}
%   Test whether a character is treated as CJK.  The ranges cover CJK
%   ideographs, compatibility and presentation forms, kana, Hangul, and the
%   punctuation LuaCJK compresses or hangs.
%    \begin{macrocode}
local function is_cjk (c)
  return (c >= 0x2E80 and c <= 0x9FFF)
      or (c >= 0xF900 and c <= 0xFAFF)
      or (c >= 0xFE30 and c <= 0xFE4F)
      or (c >= 0xFF00 and c <= 0xFFEF)
      or (c >= 0x20000 and c <= 0x3FFFF)
      or (c >= 0x3040 and c <= 0x30FF)
      or (c >= 0xAC00 and c <= 0xD7AF)
      or (c >= 0x1100 and c <= 0x11FF)
      or c == 0x2014 or c == 0x2018 or c == 0x2019
      or c == 0x201C or c == 0x201D or c == 0x2026
end

M.is_cjk = is_cjk

%    \end{macrocode}
% \end{macro}
%
% \subsection{Punctuation data}
%
% \begin{variable}[internal]{PUNCT_CLASS}
%   The punctuation LuaCJK knows about, grouped as opening, closing, middle,
%   dash, and ellipsis.  The group decides compression, line-breaking
%   prohibitions, and hanging.
%    \begin{macrocode}
local PUNCT_CLASS = {
  [0x2018] = "open",  [0x201C] = "open",
  [0x3008] = "open",  [0x300A] = "open",
  [0x300C] = "open",  [0x300E] = "open",
  [0x3010] = "open",  [0x3014] = "open",
  [0x3016] = "open",  [0x3018] = "open",
  [0x301A] = "open",  [0xFF08] = "open",
  [0xFF3B] = "open",  [0xFF5B] = "open",
  [0x2019] = "close", [0x201D] = "close",
  [0x3001] = "close", [0x3002] = "close",
  [0x3009] = "close", [0x300B] = "close",
  [0x300D] = "close", [0x300F] = "close",
  [0x3011] = "close", [0x3015] = "close",
  [0x3017] = "close", [0x3019] = "close",
  [0x301B] = "close", [0xFF01] = "close",
  [0xFF09] = "close", [0xFF0C] = "close",
  [0xFF0E] = "close", [0xFF1A] = "close",
  [0xFF1B] = "close", [0xFF1F] = "close",
  [0xFF3D] = "close", [0xFF5D] = "close",
  [0x00B7] = "middle", [0x2027] = "middle",
  [0x30FB] = "middle",
  [0x2013] = "dash",  [0x2014] = "dash",
  [0x2015] = "dash",  [0xFF0D] = "dash",
  [0x2026] = "ellipsis", [0x22EF] = "ellipsis",
}

%    \end{macrocode}
% \end{variable}
%
% \begin{macro}[internal]{punct_class}
%   Look up the group of a character.
%    \begin{macrocode}
local function punct_class (c)
  return PUNCT_CLASS[c]
end

%    \end{macrocode}
% \end{macro}
%
% \begin{variable}[internal]{KAIMING_FULL}
%   Characters that keep a full em width in the \texttt{kaiming} style.
%    \begin{macrocode}
-- Characters that keep a full em width in the kaiming style.
local KAIMING_FULL = {
  [0xFF0E] = true,  -- ．
  [0x3002] = true,  -- 。
  [0xFF1F] = true,  -- ？
  [0xFF01] = true,  -- ！
}

%    \end{macrocode}
% \end{variable}
%
% \begin{variable}[internal]{LuaCJK.punct_styles}
%   The built-in styles.  Each entry holds the compression switch, the
%   compression ratio, the opening and closing exemption, and optional fixed
%   width or hanging parameters.
%    \begin{macrocode}
M.punct_styles = {
  quanjiao = { compress_adjacent = true,  ratio = 0.5, keep_open_close = true },
  plain    = { compress_adjacent = false, ratio = 0.0, keep_open_close = true },
  banjiao  = { compress_adjacent = false, ratio = 0.0, keep_open_close = true,
               fixed_ratio = 0.5 },
  kaiming  = { compress_adjacent = false, ratio = 0.0, keep_open_close = true,
               fixed_ratio = 0.5, fixed_full = KAIMING_FULL },
  CCT      = { compress_adjacent = false, ratio = 0.0, keep_open_close = true,
               fixed_ratio = 0.7 },
  hangmobanjiao = { compress_adjacent = false, ratio = 0.0,
                    keep_open_close = true, hanging = true },
}

%    \end{macrocode}
% \end{variable}
%
% \begin{macro}[internal]{LuaCJK.define_punct_style}
%   Register a style declared by \cs{LuaCJKDeclarePunctStyle}.
%    \begin{macrocode}
function M.define_punct_style (name, compress, ratio, keep)
  M.punct_styles[name] = {
    compress_adjacent = compress,
    ratio = ratio,
    keep_open_close = keep,
  }
end

%    \end{macrocode}
% \end{macro}
%
% \begin{macro}[internal]{punct_style}
%   Read the style name from the \TeX{} token list and return the matching
%   table, falling back to \texttt{plain}.
%    \begin{macrocode}
M.punct_width = {}
M.punct_kern = {}

local function first_code (s)
  for _, code in utf8.codes(s) do
    return code
  end
end

function M.set_punct_width (chars, sp)
  for _, code in utf8.codes(chars) do
    M.punct_width[code] = sp
  end
end

function M.set_punct_kern (before, after, sp)
  local a, b = first_code(before), first_code(after)
  if a and b then
    M.punct_kern[a] = M.punct_kern[a] or {}
    M.punct_kern[a][b] = sp
  end
end

local function user_width (char)
  return M.punct_width[char]
end

local function user_kern (before, after)
  local row = M.punct_kern[before]
  return row and row[after]
end

local function punct_style ()
  local name = token.get_macro("l_LuaCJK_punct_style_tl")
  return M.punct_styles[name] or M.punct_styles.plain
end

%    \end{macrocode}
% \end{macro}
%
% \begin{macro}[internal]{kinsoku_enabled}
%   Read the line-breaking switch.
%    \begin{macrocode}
local function kinsoku_enabled ()
  return token.get_macro("l_LuaCJK_punct_kinsoku_tl") == "1"
end

%    \end{macrocode}
% \end{macro}
%
% \subsection{Suspension}
%
% \begin{variable}[internal]{LuaCJK.attr_disabled}
%   The attribute number of the suspension switch.
%    \begin{macrocode}
M.attr_disabled = M.attr_disabled or nil

%    \end{macrocode}
% \end{variable}
%
% \begin{macro}[internal]{disabled}
%   Test the suspension attribute on a node.
%    \begin{macrocode}
local function disabled (n)
  if not (M.attr_disabled and n) then
    return false
  end
  local v = has_attribute(n, M.attr_disabled)
  return v ~= nil and v ~= 0
end

%    \end{macrocode}
% \end{macro}
%
% \subsection{Font rewriting}
%
% \begin{macro}[internal]{rewrite_fonts}
%   Replace CJK glyphs by a copy in the font recorded in the attribute.
%   \begin{macro}[internal]{copy_glyph}
%     Copy a glyph node, changing only its font.
%    \begin{macrocode}
local function copy_glyph (n, fid)
  local g = new_node("glyph")
  g.font = fid
  g.char = n.char
  g.attr = n.attr
  g.lang = n.lang
  g.left = n.left
  g.right = n.right
  g.uchyph = n.uchyph
  g.xoffset = n.xoffset
  g.yoffset = n.yoffset
  g.expansion_factor = n.expansion_factor
  g.data = n.data
  g.subtype = n.subtype
  return g
end

%    \end{macrocode}
%   \end{macro}
%    \begin{macrocode}
local function rewrite_fonts (head)
  local prev = nil
  local cur = head
  while cur do
    local nxt = cur.next
    if cur.id == GLYPH and is_cjk(cur.char) and M.attr_font then
      local target = has_attribute(cur, M.attr_font)
      if target and target > 0 and target ~= cur.font then
        -- A font without the character leaves a placeholder glyph behind:
        -- luaotfload uses either the private-use range or a .notdef glyph
        -- with character 0.  Once the CJK glyph is replaced, that
        -- placeholder is no longer needed and would add a spurious advance.
        if nxt and nxt.id == GLYPH and nxt.font == cur.font
           and (nxt.char == 0
             or (nxt.char >= 0xF0000 and nxt.char <= 0x10FFFD)) then
          local after = nxt.next
          free_node(nxt)
          nxt = after
        end
        local g = copy_glyph(cur, target)
        g.next = nxt
        if prev then
          prev.next = g
        else
          head = g
        end
        free_node(cur)
        cur = g
      end
    end
    prev = cur
    cur = nxt
  end
  return head
end

%    \end{macrocode}
% \end{macro}
%
% \subsection{Node classes}
%
% \begin{macro}[internal]{class_of_char}
%   Map a character to \texttt{CJK} or \texttt{LATIN}.
%    \begin{macrocode}
local function class_of_char (c)
  if is_cjk(c) then
    return "CJK"
  end
  return "LATIN"
end

%    \end{macrocode}
% \end{macro}
%
% \begin{macro}[internal]{class_of_node}
%   Return the class at both edges of a node.  Boxes are inspected by
%   looking at their contents, and formulas count as Latin.
%   \begin{macro}[internal]{first_visible_class}
%     The class of the first visible character of a list.
%    \begin{macrocode}
local function first_visible_class (list)
  local n = list
  while n do
    if n.id == GLYPH then
      return class_of_char(n.char)
    elseif n.id == HLIST or n.id == VLIST then
      local c = first_visible_class(n.head)
      if c then
        return c
      end
    elseif n.id == MATH then
      return "LATIN"
    end
    n = n.next
  end
  return nil
end

%    \end{macrocode}
%   \end{macro}
%   \begin{macro}[internal]{last_visible_class}
%     The class of the last visible character of a list.
%    \begin{macrocode}
local function last_visible_class (list)
  local n = list
  local last = nil
  while n do
    if n.id == GLYPH then
      last = class_of_char(n.char)
    elseif n.id == HLIST or n.id == VLIST then
      local c = last_visible_class(n.head)
      if c then
        last = c
      end
    elseif n.id == MATH then
      last = "LATIN"
    end
    n = n.next
  end
  return last
end

%    \end{macrocode}
%   \end{macro}
%    \begin{macrocode}
local function class_of_node (n)
  if n.id == GLYPH then
    local c = class_of_char(n.char)
    return c, c
  elseif n.id == HLIST or n.id == VLIST then
    local first = first_visible_class(n.head)
    local last = last_visible_class(n.head)
    if first and last then
      return first, last
    end
  elseif n.id == MATH then
    return "LATIN", "LATIN"
  end
  return nil, nil
end

%    \end{macrocode}
% \end{macro}
%
% \subsection{Glue and kern helpers}
%
% \begin{macro}[internal]{is_inserted}
%   Test whether a node was inserted by LuaCJK.
%    \begin{macrocode}
local function is_inserted (n)
  return M.attr_inserted
    and has_attribute(n, M.attr_inserted) ~= nil
end

%    \end{macrocode}
% \end{macro}
%
% \begin{macro}[internal]{get_glue}
%   Read a skip register by name.
%    \begin{macrocode}
local function get_glue (name)
  local s = tex.skip[name]
  if not s then
    return nil
  end
  return s
end

%    \end{macrocode}
% \end{macro}
%
% \begin{macro}[internal]{inherit_attributes}
%   Give a new node the attribute list of the node it is spliced next to.
%   LuaCJK creates glue, kerns, and penalties after \TeX{} has attached the
%   surrounding attributes to the text, so a package that marks spans with
%   attributes, such as \pkg{lua-ul}, would otherwise read the new node as
%   the end of its span and start a new one.  Assigning to the
%   \meta{attr} field is reference counted, and later per-attribute changes
%   are copy-on-write, so the source node is never modified.
%    \begin{macrocode}
local function inherit_attributes (new, source)
  if source and source.attr then
    new.attr = source.attr
  end
end

%    \end{macrocode}
% \end{macro}
%
% \begin{macro}[internal]{make_glue}
%   Create a glue node from a skip register, inherit the attributes of
%   \meta{source}, and mark the node as inserted.
%    \begin{macrocode}
local function make_glue (name, source)
  local s = get_glue(name)
  if not s then
    return nil
  end
  local g = new_node("glue")
  inherit_attributes(g, source)
  g.width = s.width
  g.stretch = s.stretch
  g.shrink = s.shrink
  g.stretch_order = 0
  g.shrink_order = 0
  if M.attr_inserted then
    set_attribute(g, M.attr_inserted, 1)
  end
  return g
end

%    \end{macrocode}
% \end{macro}
%
% \begin{macro}[internal]{set_glue}
%   Give an existing glue node the dimensions of a skip register.
%    \begin{macrocode}
local function set_glue (g, name)
  local s = get_glue(name)
  if not s then
    return false
  end
  g.width = s.width
  g.stretch = s.stretch
  g.shrink = s.shrink
  if M.attr_inserted then
    set_attribute(g, M.attr_inserted, 1)
  end
  return true
end

%    \end{macrocode}
% \end{macro}
%
% \begin{macro}[internal]{make_kern}
%   Create a kern node, inherit the attributes of \meta{source}, and mark
%   the node as inserted.
%    \begin{macrocode}
local function make_kern (width, source)
  local k = new_node("kern")
  inherit_attributes(k, source)
  k.kern = width
  if M.attr_inserted then
    set_attribute(k, M.attr_inserted, 1)
  end
  return k
end

%    \end{macrocode}
% \end{macro}
%
% \subsection{Spacing}
%
% \begin{macro}[internal]{spacing}
%   Walk a list and insert the glue between CJK and Latin material.  The
%   walk tracks the last visible node and the class at its end, so that a
%   formula is treated as one Latin atom.
%   \begin{macro}[internal]{handle_boundary}
%     Decide what to do at the boundary between two visible nodes.  Only
%     source-generated glue is a candidate for replacement; other glue and
%     penalties are left alone.
%     \begin{macro}[internal]{insert_boundary_glue}
%       Insert the CJK--CJK or CJK--Latin glue, or the compression kern for
%       two adjacent punctuation marks.
%      \begin{macrocode}
M.normal_spaced = {}

function M.set_normal_spaced_chars (chars)
  local set = {}
  for _, code in utf8.codes(chars) do
    set[code] = true
  end
  M.normal_spaced = set
end

local function normal_spaced (n)
  return n and n.id == GLYPH and M.normal_spaced[n.char] == true
end

local function insert_boundary_glue (head, prev, prev_end, cur, cur_start)
  local name = nil
  if prev_end == "CJK" and cur_start == "CJK" then
    name = "l_LuaCJK_ccglue_skip"
  elseif prev_end ~= cur_start
     and not normal_spaced(prev) and not normal_spaced(cur) then
    name = "l_LuaCJK_ecglue_skip"
  end

  local style = punct_style()
  local prev_punct = prev and prev.id == GLYPH and punct_class(prev.char)
  local cur_punct = cur.id == GLYPH and punct_class(cur.char)
  if prev_punct and cur_punct then
    local explicit = user_kern(prev.char, cur.char)
    if explicit then
      local k = make_kern(explicit, prev)
      k.next = cur
      if prev then
        prev.next = k
      else
        head = k
      end
      return head
    end
  end
  if prev_punct and cur_punct and style.compress_adjacent then
    local exempt = (style.keep_open_close
        and prev_punct == "open" and cur_punct == "close")
      or (prev_punct == cur_punct
        and (prev_punct == "dash" or prev_punct == "ellipsis"))
    if not exempt and prev.width and prev.width > 0 then
      local k = make_kern(-math.floor(style.ratio * prev.width + 0.5), prev)
      local g = name and make_glue(name, prev) or nil
      k.next = g or cur
      if g then
        g.next = cur
      end
      if prev then
        prev.next = k
      else
        head = k
      end
      return head
    end
  end

  if not name then
    return head
  end
  local g = make_glue(name, prev or cur)
  if not g then
    return head
  end
  g.next = cur
  if prev then
    prev.next = g
  else
    head = g
  end
  return head
end

%    \end{macrocode}
%     \end{macro}
%     \begin{macro}[internal]{handle_source_space}
%       Remove or retime a space typed in the source.
%      \begin{macrocode}
local function handle_source_space (head, glue, before, prev_end, cur, cur_start)
  -- Only source-generated spaces are candidates.  Explicit \hskip, \hfill
  -- and similar glue have a different subtype and must be preserved.
  if glue.subtype ~= SPACESKIP then
    return head
  end
  local prev_cjk = prev_end == "CJK"
  local cur_cjk = cur_start == "CJK"
  if prev_cjk and cur_cjk then
    local cjk_space = M.attr_cjk_space
      and (has_attribute(glue, M.attr_cjk_space) or 0) ~= 0
    if not cjk_space then
      if before then
        before.next = cur
      else
        head = cur
      end
      free_node(glue)
    end
  elseif prev_cjk or cur_cjk then
    local xecglue = M.attr_xecglue
      and (has_attribute(glue, M.attr_xecglue) or 0) ~= 0
    if xecglue then
      set_glue(glue, "l_LuaCJK_ecglue_skip")
    end
  end
  return head
end

%    \end{macrocode}
%     \end{macro}
%    \begin{macrocode}
local function handle_boundary (head, prev, before, prev_visible,
                                prev_end, cur, start_class)
  if not prev_visible then
    return head
  end
  if disabled(prev_visible) or disabled(cur) then
    return head
  end
  if prev and prev.id == GLUE then
    if not is_inserted(prev) then
      return handle_source_space(head, prev, before, prev_end, cur, start_class)
    end
  elseif prev and prev.id == PENALTY then
    -- A penalty is an explicit barrier: do not bridge it.
  elseif prev and prev.id == DISC then
    -- Discretionaries are opaque in this first boundary version.
  else
    return insert_boundary_glue(head, prev, prev_end, cur, start_class)
  end
  return head
end

%    \end{macrocode}
%   \end{macro}
%    \begin{macrocode}
local function spacing (head)
  local prev = nil
  local before = nil
  local prev_visible = nil
  local prev_end = nil
  local math_level = 0
  local cur = head
  while cur do
    local nxt = cur.next
    if cur.id == MATH then
      if cur.subtype == 0 then
        if math_level == 0 then
          -- The formula behaves like a Latin atom at its left edge.
          head = handle_boundary(
            head, prev, before, prev_visible, prev_end, cur, "LATIN")
        end
        math_level = math_level + 1
        prev_visible = nil
        prev_end = nil
      else
        if math_level > 0 then
          math_level = math_level - 1
        end
        if math_level == 0 then
          prev_visible = cur
          prev_end = "LATIN"
        end
      end
    elseif math_level > 0 then
      -- Inside a formula the boundary rules do not apply.
    else
      local start_class, end_class = class_of_node(cur)
      if start_class then
        head = handle_boundary(
          head, prev, before, prev_visible, prev_end, cur, start_class)
        prev_visible = cur
        prev_end = end_class
      elseif cur.id == PENALTY or cur.id == DISC then
        prev_visible = nil
        prev_end = nil
      end
    end
    before = prev
    prev = cur
    cur = nxt
  end
  return head
end

%    \end{macrocode}
% \end{macro}
%
% \subsection{Line-breaking prohibitions}
%
% \TeX{} may break at glue and at kern.  Both are places where LuaCJK inserts
% its own nodes, so after building the list a second pass puts a penalty
% 10000 in front of every break opportunity that must not be used.
%
% \begin{macro}[internal]{apply_kinsoku}
%   Add the prohibitions.
%    \begin{macrocode}
-- Line-breaking prohibitions (kinsoku).
--
-- TeX may break at glue and at kern.  Both are the places we insert our own
-- spacing nodes, so after building the list we walk it once more and put a
-- penalty 10000 in front of every break opportunity that must not be used.
%    \end{macrocode}
%   \begin{macro}[internal]{glyph_class}
%     The punctuation group of a glyph.
%    \begin{macrocode}
local function glyph_class (n)
  if n and n.id == GLYPH then
    return punct_class(n.char)
  end
  return nil
end

%    \end{macrocode}
%   \end{macro}
%   \begin{macro}[internal]{next_visible}
%     The next visible glyph, stopping at an explicit break barrier.
%    \begin{macrocode}
local function next_visible (n)
  while n do
    if n.id == GLYPH then
      return n
    elseif n.id == HLIST or n.id == VLIST then
      if first_visible_class(n.head) then
        return nil
      end
    elseif n.id == PENALTY or n.id == DISC then
      return nil
    end
    n = n.next
  end
  return nil
end

%    \end{macrocode}
%   \end{macro}
%   \begin{macro}[internal]{forbids_break}
%     The rules: closing and middle punctuation must not start a line,
%     opening punctuation must not end one, and a run of dashes or ellipses
%     stays together.
%    \begin{macrocode}
local function forbids_break (prev_punct, next_punct)
  -- A closing or middle punctuation must not start a line.
  if next_punct == "close" or next_punct == "middle" then
    return true
  end
  -- An opening punctuation must not end a line.
  if prev_punct == "open" then
    return true
  end
  -- Long punctuation is one unit: no break inside a run.
  if prev_punct == next_punct
     and (prev_punct == "dash" or prev_punct == "ellipsis") then
    return true
  end
  return false
end

%    \end{macrocode}
%   \end{macro}
%   \begin{macro}[internal]{is_break_point}
%     Glue and kerns are break opportunities; only the ones LuaCJK inserted
%     and the source spaces are considered.
%    \begin{macrocode}
local function is_break_point (n)
  if n.id == GLUE then
    return is_inserted(n) or n.subtype == SPACESKIP
  elseif n.id == KERN then
    return is_inserted(n)
  end
  return false
end

%    \end{macrocode}
%   \end{macro}
%    \begin{macrocode}
local function apply_kinsoku (head)
  if not kinsoku_enabled() then
    return head
  end
  local prev = nil
  local last_visible = nil
  local cur = head
  while cur do
    if cur.id == GLYPH then
      last_visible = cur
    elseif is_break_point(cur) and last_visible then
      local nxt = next_visible(cur.next)
      if nxt and not disabled(last_visible) and not disabled(nxt)
         and forbids_break(glyph_class(last_visible), glyph_class(nxt)) then
        if not (prev and prev.id == PENALTY and prev.penalty >= 10000) then
          local p = new_node("penalty")
          inherit_attributes(p, last_visible)
          p.penalty = 10000
          if M.attr_inserted then
            set_attribute(p, M.attr_inserted, 1)
          end
          p.next = cur
          if prev then
            prev.next = p
          else
            head = p
          end
          cur = p
        end
      end
    end
    prev = cur
    cur = cur.next
  end
  return head
end

%    \end{macrocode}
% \end{macro}
%
% \begin{macro}[internal]{apply_fixed_widths}
%   Styles such as \texttt{banjiao} give every punctuation mark a fixed
%   advance.  The glyph advance cannot be changed, so a compensating kern
%   goes right after the glyph.  A penalty in front of that kern keeps \TeX{}
%   from treating it as an extra break opportunity.
%    \begin{macrocode}
-- Fixed punctuation width.
--
-- Styles such as banjiao give every punctuation mark a fixed advance.  The
-- glyph advance cannot be changed, so a compensating kern goes right after
-- the glyph.  A penalty in front of that kern keeps TeX from treating it as
-- an extra break opportunity.
local function apply_fixed_widths (head)
  local style = punct_style()
  local ratio = style.fixed_ratio
  local full = style.fixed_full
  local cur = head
  while cur do
    if cur.id == GLYPH and punct_class(cur.char) and not disabled(cur) then
      local target = user_width(cur.char)
      if not target and ratio then
        local f = font.getfont(cur.font)
        local size = f and f.size or 0
        if size > 0 and not (full and full[cur.char]) then
          target = math.floor(ratio * size + 0.5)
        end
      end
      if target then
        local delta = target - cur.width
        if delta ~= 0 then
          local p = new_node("penalty")
          inherit_attributes(p, cur)
          p.penalty = 10000
          local k = make_kern(delta, cur)
          p.next = k
          k.next = cur.next
          cur.next = p
        end
      end
    end
    cur = cur.next
  end
  return head
end

%    \end{macrocode}
% \end{macro}
%
% \subsection{CJK in mathematics}
%
% A CJK character in math mode reaches the math list as a \texttt{math_char}
% noad.  Its family points at the math font, which usually has no CJK glyph,
% so LuaCJK replaces the nucleus by a \texttt{sub_box} holding the character
% in the CJK font.
%
% \begin{variable}[internal]{MATH_CHAR, MATH_TEXT_CHAR}
%   The node types of the two kinds of math character noad.
%    \begin{macrocode}
-- CJK in math mode.
--
-- A CJK character in math mode reaches the math list as a `math_char` noad.
-- Its family points at the math font, which usually has no CJK glyph.  We
-- replace the nucleus by a `sub_box` holding the character in the CJK font.
local MATH_CHAR = node.id("math_char")
local MATH_TEXT_CHAR = node.id("math_text_char")

%    \end{macrocode}
% \end{variable}
%
% \begin{macro}[internal]{math_font_id}
%   Read the font identifier chosen by \TeX.
%    \begin{macrocode}
local function math_font_id ()
  local fid = tex.count["g_LuaCJK_math_fontid_int"]
  if fid and fid > 0 then
    return fid
  end
  return nil
end

%    \end{macrocode}
% \end{macro}
%
% \begin{macro}[internal]{cjk_math_enabled}
%   Read the switch set by the \texttt{CJKmath} option.
%    \begin{macrocode}
local function cjk_math_enabled ()
  return token.get_macro("l_LuaCJK_cjk_math_tl") == "1"
end

%    \end{macrocode}
% \end{macro}
%
% \begin{variable}[internal]{SUB_MLIST}
%   The node type of a nested math list.
%    \begin{macrocode}
local SUB_MLIST = node.id("sub_mlist")

%    \end{macrocode}
% \end{variable}
%
% \begin{macro}[internal]{rewrite_math}
%   Rewrite the CJK noads of a math list, or warn once when the feature is
%   off.
%   \begin{macro}[internal]{rewrite_math_list}
%     Walk a math list.  The declaration is split from the definition
%     because the walk recurses into sub lists.
%    \begin{macrocode}
-- Recurse into sub lists (fractions, braces, scripts) so that nested CJK
-- characters are rewritten too.
local rewrite_math_list

%    \end{macrocode}
%     \begin{macro}[internal]{cjk_noad}
%       Build a noad whose nucleus is a box with the character in the CJK
%       font.
%      \begin{macrocode}
local function cjk_noad (char, fid)
  local g = new_node("glyph")
  g.font = fid
  g.char = char
  local b = new_node("hlist")
  b.head = g
  b.width = g.width
  b.height = g.height
  b.depth = g.depth
  local sb = new_node("sub_box")
  sb.head = b
  local n = new_node("noad")
  n.nucleus = sb
  return n
end

%    \end{macrocode}
%     \end{macro}
%     \begin{macro}[internal]{rewrite_noad}
%       Replace the nucleus of one noad and recurse into the sub lists of its
%       numerator, denominator, superscript, and subscript.
%      \begin{macrocode}
local function rewrite_noad (noad, fid)
  if disabled(noad) then
    return
  end
  local nuc = noad.nucleus
  if nuc then
    if nuc.id == MATH_CHAR or nuc.id == MATH_TEXT_CHAR then
      if is_cjk(nuc.char) then
        if M.attr_math_font then
          local a = has_attribute(nuc, M.attr_math_font)
          if a and a > 0 then
            fid = a
          end
        end
        noad.nucleus = cjk_noad(nuc.char, fid).nucleus
      end
    elseif nuc.id == SUB_MLIST then
      rewrite_math_list(nuc.head, fid)
    end
  end
  for _, field in ipairs({ "num", "denom", "sup", "sub" }) do
    local s = noad[field]
    if s and s.id == SUB_MLIST then
      rewrite_math_list(s.head, fid)
    end
  end
end

%    \end{macrocode}
%     \end{macro}
%    \begin{macrocode}
rewrite_math_list = function (head, fid)
  local cur = head
  while cur do
    rewrite_noad(cur, fid)
    cur = cur.next
  end
  return head
end

%    \end{macrocode}
%   \end{macro}
%
%   \begin{variable}[internal]{LuaCJK.warned_cjk_math}
%     Whether the one-time warning has been shown.
%    \begin{macrocode}
M.warned_cjk_math = M.warned_cjk_math or false

%    \end{macrocode}
%   \end{variable}
%
%   \begin{macro}[internal]{warn_disabled_cjk_math}
%     Warn once when a CJK character appears in math mode while the feature
%     is off.
%     \begin{macro}[internal]{find_cjk_math}
%       Search a math list, including its sub lists, for a CJK character.
%      \begin{macrocode}
local function find_cjk_math (head)
  local cur = head
  while cur do
    local nuc = cur.nucleus
    if nuc then
      if (nuc.id == MATH_CHAR or nuc.id == MATH_TEXT_CHAR)
         and is_cjk(nuc.char) then
        return true
      elseif nuc.id == SUB_MLIST and find_cjk_math(nuc.head) then
        return true
      end
    end
    for _, field in ipairs({ "num", "denom", "sup", "sub" }) do
      local sub = cur[field]
      if sub and sub.id == SUB_MLIST and find_cjk_math(sub.head) then
        return true
      end
    end
    cur = cur.next
  end
  return false
end

%    \end{macrocode}
%     \end{macro}
%    \begin{macrocode}
local function warn_disabled_cjk_math (head)
  if M.warned_cjk_math then
    return
  end
  if find_cjk_math(head) then
    M.warned_cjk_math = true
    texio.write_nl("term and log",
      "Package LuaCJK Warning: CJK character in math mode, but CJKmath is not enabled.")
    texio.write_nl("term and log",
      "(LuaCJK)                Use \\LuaCJKsetup{CJKmath=true} " ..
      "to typeset it with the CJK font.")
  end
end

%    \end{macrocode}
%   \end{macro}
%    \begin{macrocode}
local function rewrite_math (head)
  if not cjk_math_enabled() then
    warn_disabled_cjk_math(head)
    return head
  end
  local fid = math_font_id()
  if not fid then
    return head
  end
  return rewrite_math_list(head, fid)
end

%    \end{macrocode}
% \end{macro}
%
% \subsection{Hanging punctuation}
%
% The glyph advance stays inside the line, so shifting the ink sideways lets
% a line-final closing punctuation hang into the margin without changing the
% line width.  Line-initial opening punctuation is shifted the other way.
%
% \begin{macro}[internal]{hang_punctuation}
%   Apply the hanging shift to every line.
%   \begin{macro}[internal]{hang_line}
%     Shift the first and last punctuation glyph of one line.
%    \begin{macrocode}
-- Hanging punctuation (hangmobanjiao).
--
-- The glyph advance stays inside the line, so shifting the ink sideways lets
-- a line-final closing punctuation hang into the margin without changing the
-- line width.  Line-initial opening punctuation is shifted the other way.
%    \end{macrocode}
%     \begin{macro}[internal]{find_first_last_glyph}
%       The first and last glyph nodes of a list.
%      \begin{macrocode}
local function find_first_last_glyph (list)
  local first, last
  local n = list
  while n do
    if n.id == GLYPH then
      if not first then first = n end
      last = n
    end
    n = n.next
  end
  return first, last
end

%    \end{macrocode}
%     \end{macro}
%     \begin{macro}[internal]{shift_glyph}
%       Move the ink of a glyph by a horizontal offset.
%      \begin{macrocode}
local function shift_glyph (g, delta)
  if delta == 0 then return end
  g.xoffset = g.xoffset + delta
end

%    \end{macrocode}
%     \end{macro}
%    \begin{macrocode}
local function hang_line (line)
  if line.id ~= HLIST then
    return
  end
  local first, last = find_first_last_glyph(line.head)
  if last and disabled(last) then
    last = nil
  end
  if first and disabled(first) then
    first = nil
  end
  if last and punct_class(last.char) then
    local cls = punct_class(last.char)
    if cls == "close" or cls == "middle" then
      local f = font.getfont(last.font)
      local size = f and f.size or 0
      shift_glyph(last, math.floor(0.5 * size + 0.5))
    end
  end
  if first and punct_class(first.char) then
    local cls = punct_class(first.char)
    if cls == "open" then
      local f = font.getfont(first.font)
      local size = f and f.size or 0
      shift_glyph(first, -math.floor(0.5 * size + 0.5))
    end
  end
end

%    \end{macrocode}
%   \end{macro}
%    \begin{macrocode}
local function hang_punctuation (head)
  local style = punct_style()
  if not style.hanging then
    return head
  end
  local n = head
  while n do
    hang_line(n)
    n = n.next
  end
  return head
end

%    \end{macrocode}
% \end{macro}
%
% \subsection{Callbacks}
%
% \begin{macro}[internal]{LuaCJK.process}
%   Run the font, spacing, fixed-width, and kinsoku passes over a list.
%    \begin{macrocode}
local function process (head)
  head = rewrite_fonts(head)
  head = spacing(head)
  head = apply_fixed_widths(head)
  head = apply_kinsoku(head)
  return head
end

M.process = process

%    \end{macrocode}
% \end{macro}
%
% \begin{macro}[internal]{LuaCJK.pre_linebreak}
%   The line-breaking callback.
%    \begin{macrocode}
function M.pre_linebreak (head, groupcode)
  return process(head)
end

%    \end{macrocode}
% \end{macro}
%
% \begin{macro}[internal]{LuaCJK.hpack}
%   The box-packing callback.
%    \begin{macrocode}
function M.hpack (head, groupcode, size, packtype, direction, attributelist)
  return process(head)
end

%    \end{macrocode}
% \end{macro}
%
% \begin{macro}[internal]{LuaCJK.pre_mlist_to_hlist}
%   The math callback.
%    \begin{macrocode}
function M.pre_mlist_to_hlist (head, display, penalties)
  return rewrite_math(head)
end

%    \end{macrocode}
% \end{macro}
%
% \begin{macro}[internal]{LuaCJK.post_linebreak}
%   The hanging punctuation callback.
%    \begin{macrocode}
function M.post_linebreak (head, groupcode)
  return hang_punctuation(head)
end

%    \end{macrocode}
% \end{macro}
%
% The callbacks are registered in this order: spacing, font correction,
% math, and hanging.
%    \begin{macrocode}
luatexbase.add_to_callback("pre_linebreak_filter", M.pre_linebreak, "LuaCJK.spacing")
luatexbase.add_to_callback("hpack_filter", M.hpack, "LuaCJK.spacing")
luatexbase.add_to_callback("pre_mlist_to_hlist_filter", M.pre_mlist_to_hlist, "LuaCJK.math")
luatexbase.add_to_callback("post_linebreak_filter", M.post_linebreak, "LuaCJK.hang")
%</lua-core>
%    \end{macrocode}
%

% \subsection{Font effects}
%
% The \pkg*{LuaCJKfntef} package is a thin layer over \pkg{lua-ul}, which
% provides the underlining engine.  It is generated from this file and may be
% loaded on its own; it loads LuaCJK first.
%    \begin{macrocode}
%<*fntef>
\sys_if_engine_luatex:F
  {
    \msg_new:nnn { LuaCJKfntef } { engine }
      { LuaCJKfntef~requires~LuaTeX. }
    \msg_critical:nn { LuaCJKfntef } { engine }
    \tex_endinput:D
  }
\RequirePackage { LuaCJK }
\RequirePackage { lua-ul }
%    \end{macrocode}
%
% \begin{macro}{\CJKunderline}
%   Underline the argument; the optional argument is passed to
%   \pkg{lua-ul}.
%    \begin{macrocode}
\NewDocumentCommand \CJKunderline { s t- o +m }
  {
    \IfValueTF {#3}
      { \underLine [ #3 ] { #4 } }
      { \underLine { #4 } }
  }
%    \end{macrocode}
% \end{macro}
%
% \begin{macro}{\CJKsout}
%   Strike through the argument.
%    \begin{macrocode}
\NewDocumentCommand \CJKsout { s t- o +m }
  {
    \IfValueTF {#3}
      { \strikeThrough [ #3 ] { #4 } }
      { \strikeThrough { #4 } }
  }
%    \end{macrocode}
% \end{macro}
%
% \begin{macro}{\CJKunderdot}
%   Emphasise with dots below the glyphs.  The optional argument of
%   \pkg{xeCJK} is not implemented, so it only raises a warning.
%    \begin{macrocode}
\NewDocumentCommand \CJKunderdot { s t- o +m }
  {
    \IfValueT {#3}
      {
        \msg_warning:nn { LuaCJKfntef } { no-options }
      }
    \LuaCJKfntef_underdot:n { #4 }
  }
\msg_new:nnn { LuaCJKfntef } { no-options }
  { The~optional~argument~of~`\token_to_str:N \CJKunderdot'~is~ignored. }
%    \end{macrocode}
%   \begin{macro}[internal]{\LuaCJKfntef_underdot:n}
%     Apply the dot pattern inside a group.
%     \begin{variable}[internal]
%       {\l_LuaCJKfntef_dot_dim, \l_LuaCJKfntef_dot_gap_dim,
%        \l_LuaCJKfntef_dot_raise_dim, \l_LuaCJKfntef_dots_box}
%       The size of one dot, the space on both sides of it, the amount it is
%       lowered, and the box holding one repeating unit of the pattern.
%      \begin{macrocode}
\dim_new:N \l_LuaCJKfntef_dot_dim
\dim_set:Nn \l_LuaCJKfntef_dot_dim { .06em }
\dim_new:N \l_LuaCJKfntef_dot_gap_dim
\dim_set:Nn \l_LuaCJKfntef_dot_gap_dim { .12em }
\dim_new:N \l_LuaCJKfntef_dot_raise_dim
\dim_set:Nn \l_LuaCJKfntef_dot_raise_dim { -.40em }
\box_new:N \l_LuaCJKfntef_dots_box
\hbox_set:Nn \l_LuaCJKfntef_dots_box
  {
    \skip_horizontal:n { \l_LuaCJKfntef_dot_gap_dim }
    \raise \l_LuaCJKfntef_dot_raise_dim \tex_hbox:D
      {
        \vrule width \l_LuaCJKfntef_dot_dim
               height \l_LuaCJKfntef_dot_dim
               depth \c_zero_dim
      }
    \skip_horizontal:n { \l_LuaCJKfntef_dot_gap_dim }
  }
%    \end{macrocode}
%     \end{variable}
%     \begin{macro}[internal]{\LuaCJKfntef_dots}
%       The underline type that repeats the dot box.
%      \begin{macrocode}
\newunderlinetype \LuaCJKfntef_dots
  { \leaders \box_use_drop:N \l_LuaCJKfntef_dots_box }
%    \end{macrocode}
%     \end{macro}
%     \begin{macrocode}
\cs_new_protected:Npn \LuaCJKfntef_underdot:n #1
  {
    \group_begin:
      \LuaCJKfntef_dots
      #1
    \group_end:
  }
%</fntef>
%    \end{macrocode}
%   \end{macro}
% \end{macro}
%
% \end{implementation}
%
% \Finale
%
\endinput
